\chapter{Metric Spaces and Completeness}\label{ch:metric}

\begin{center}\itshape
Press cosine often enough and the display settles on $0.739085\ldots$, from any start. Shrink the line twice by a factor $3$ and the Cantor set appears, from any start. Why does the iteration converge, and why to only one limit?
\end{center}

Set a calculator to radians, type a number, and press cosine again and again. After a few dozen presses the display freezes on $0.7390851\ldots$, and the starting number does not matter. A professor of French named Dottie is said to have found the number by pressing keys with nothing particular in mind, and the number now carries her name. The reason is short. On $[0,1]$, cosine pulls points together by a fixed ratio $k<1$. A map that does that, on a complete space, has exactly one point that it leaves still. The iterates walk to that point.

The only ingredient is a distance. Maurice Fr\'echet saw this in his thesis of 1906: functions, curves and sequences become points of a space as soon as one says how far apart they are. Felix Hausdorff gave these spaces the name \emph{metric spaces} in 1914.

\portrait{frechet}{Maurice Fr\'echet (1878--1973). Photograph, author unknown. Public domain, via Wikimedia Commons.}

The property that matters is \emph{completeness}: every Cauchy sequence converges, so the space has no holes. Chapter~\ref{ch:sets} completed $\Q$ into $\R$. The same construction completes an arbitrary metric space. Run it on $\Q$ with a stranger distance, one in which a number is small when a high power of a prime $p$ divides it, and a second field appears, $\Q_p$.

The hero of the chapter is Stefan Banach's fixed point theorem (1922). It solves equations. It proves that differential equations have solutions. It proves that Newton's method converges once the map is a contraction. It ranks the pages of the web. It draws fractals: the Cantor set and the Sierpi\'nski triangle are fixed points of maps that act on shapes. Along the way, $\exp$, $\cos$, $\sin$ and $\pi$, used since Volume~1, are defined by power series. The number $\pi$ appears as twice the first positive zero of the cosine, not as a length read off a picture of a circle.

\section{Distances}

\begin{definition}
A \emph{metric} on a set $X$ is a map $d:X\times X\to[0,+\infty[$ such that $d(x,y)=0$ if and only if $x=y$, such that $d(x,y)=d(y,x)$, and such that the triangle inequality $d(x,z)\le d(x,y)+d(y,z)$ holds. A \emph{norm} on a real or complex vector space $E$ is a map $\|\cdot\|:E\to[0,+\infty[$ with $\|x\|=0$ if and only if $x=0$, with $\|\lambda x\|=|\lambda|\|x\|$, and with $\|x+y\|\le\|x\|+\|y\|$. Then $d(x,y)=\|x-y\|$ is a metric.
\end{definition}

\begin{example}\label{ex:metrics}
\begin{enumerate}[label=(\alph*)]
\item On $\R^n$ the three norms are $\|x\|_1=\sum|x_i|$, $\|x\|_2=\bigl(\sum x_i^2\bigr)^{1/2}$ and $\|x\|_\infty=\max|x_i|$. The triangle inequality for $\|\cdot\|_2$ is Minkowski's inequality, which follows from Cauchy--Schwarz: $\|x+y\|_2^2=\|x\|_2^2+2\langle x,y\rangle+\|y\|_2^2\le(\|x\|_2+\|y\|_2)^2$. The three norms compare by $\|x\|_\infty\le\|x\|_2\le\|x\|_1\le n\|x\|_\infty$.
\item On the set $B(S)$ of bounded functions $S\to\R$ or $S\to\C$, the \emph{sup norm} is $\|f\|_\infty=\sup_{s\in S}|f(s)|$.
\item On an arbitrary set the \emph{discrete metric} is $d(x,y)=1$ when $x\ne y$. On words of length $n$, the \emph{Hamming distance} counts the positions where two words differ. That is the distance of error-correcting codes.
\item On $\Q$, for a prime $p$, the \emph{$p$-adic absolute value} is $|x|_p=p^{-v_p(x)}$, with $|0|_p=0$. The exponent $v_p(x)$ is the power of $p$ in $x\ne0$: write $x=p^v\frac ab$ with $a$ and $b$ integers not divisible by $p$, and set $v_p(x)=v$. For an integer this is the valuation of Volume~2, Chapter~1. For example $|50|_5=\frac1{25}$, $|\frac3{10}|_5=5$ and $|7|_5=1$. It is multiplicative, $|xy|_p=|x|_p|y|_p$, and it satisfies the \emph{ultrametric} inequality $|x+y|_p\le\max(|x|_p,|y|_p)$, which is stronger than the triangle inequality, because $v_p(x+y)\ge\min(v_p(x),v_p(y))$. So $d(x,y)=|x-y|_p$ is a metric, and in this metric $p^n\to0$.
\end{enumerate}
\end{example}

\begin{figure}[ht]
\centering
\fig{ch02-balls}
\caption{The unit balls of $\|\cdot\|_1$ (a diamond), $\|\cdot\|_2$ (a disc) and $\|\cdot\|_\infty$ (a square) in $\R^2$. They nest exactly as the inequalities of Example~\ref{ex:metrics}(a) say. The three norms have the same convergent sequences.}
\label{fig:balls}
\end{figure}

A subset $Y$ of a metric space is itself a metric space, with the restricted metric. The \emph{open ball} is $B(x,r)=\{y:d(x,y)<r\}$, the \emph{closed ball} is $\bar B(x,r)=\{y:d(x,y)\le r\}$, and the \emph{diameter} of a set $A$ is $\operatorname{diam}A=\sup_{x,y\in A}d(x,y)$. The \emph{distance} from a point $x$ to a non-empty set $A$ is $d(x,A)=\inf_{a\in A}d(x,a)$. It is $1$-Lipschitz in $x$: $d(x,A)\le d(x,a)\le d(x,y)+d(y,a)$, and the infimum over $a$ keeps the inequality.

\begin{figure}[ht]
\centering
\fig{ch02-rw-open}
\caption{An open ball does not contain its boundary. A sequence in the ball can approach the boundary point and have no limit inside the ball. That is why $]0,1]$ is not complete.}
\label{fig:rwopen}
\end{figure}

\section{Open sets, convergence and continuity}

\begin{definition}
A set $U\subseteq X$ is \emph{open} when every $x\in U$ has some $r>0$ with $B(x,r)\subseteq U$. It is \emph{closed} when $X\setminus U$ is open. The \emph{interior} $\mathring A$ is the union of the open sets contained in $A$. The \emph{closure} $\bar A$ is the intersection of the closed sets that contain $A$. The set $A$ is \emph{dense} when $\bar A=X$. A sequence satisfies $x_n\to x$ when $d(x_n,x)\to0$.
\end{definition}

An open ball is an open set. If $y\in B(x,r)$, the triangle inequality puts $B(y,r-d(x,y))$ inside $B(x,r)$. Arbitrary unions of open sets are open, and so are finite intersections: take the smallest radius. Intersections of closed sets are closed, and so are finite unions. A limit is unique, because $d(x,y)\le d(x,x_n)+d(x_n,y)$.

\begin{proposition}\label{prop:closure}
\begin{enumerate}[label=(\alph*)]
\item A point $x$ lies in $\bar A$ if and only if $d(x,A)=0$, and if and only if some sequence of points of $A$ converges to $x$.
\item The set $A$ is closed if and only if it contains every limit of a convergent sequence drawn from $A$.
\item For a map $f:X\to Y$ between metric spaces, the following are equivalent. (i) For every $x$ and every $\varepsilon>0$ there is a $\delta>0$ such that $d(x,y)<\delta$ implies $d(f(x),f(y))<\varepsilon$ (\emph{continuity}). (ii) $x_n\to x$ implies $f(x_n)\to f(x)$. (iii) The preimage $f^{-1}(V)$ is open whenever $V\subseteq Y$ is open.
\end{enumerate}
\end{proposition}
\begin{proof}
(a) If $d(x,A)>0$, the ball $B(x,d(x,A))$ is open and misses $A$, so its complement is a closed set that contains $A$ and not $x$. Thus $x\notin\bar A$. If $d(x,A)=0$, choose $a_n\in A$ with $d(x,a_n)<\frac1n$. Then $a_n\to x$. If $a_n\to x$ with $a_n\in A$, and if $F\supseteq A$ is closed with $x\notin F$, some ball $B(x,r)$ misses $F$ and therefore misses every $a_n$, which is impossible for large $n$. So $x\in\bar A$. (b) The set $A$ is closed if and only if $A=\bar A$. Apply (a). (c) The implication (i)$\Rightarrow$(ii) is the definition read along a sequence. For (ii)$\Rightarrow$(iii), if $f^{-1}(V)$ failed to be open there would exist $x\in f^{-1}(V)$ and a sequence $x_n\to x$ with $f(x_n)\notin V$. Then $f(x_n)\to f(x)\in V$. But $V$ contains a ball about $f(x)$, and $f(x_n)$ would have to enter that ball. For (iii)$\Rightarrow$(i), the preimage $f^{-1}(B(f(x),\varepsilon))$ is open and contains $x$, so it contains some ball $B(x,\delta)$.
\end{proof}

By (c), continuity depends only on the open sets, not on the particular metric. That observation is the start of Chapter~\ref{ch:topology}. Two metrics are \emph{equivalent} when they determine the same open sets, and \emph{Lipschitz equivalent} when $c\,d\le d'\le C\,d$ for constants $0<c\le C$. The three norms of Example~\ref{ex:metrics}(a) are Lipschitz equivalent. Chapter~\ref{ch:compact} proves that every pair of norms on $\R^n$ is.

\begin{definition}
A map $f:X\to Y$ is \emph{uniformly continuous} when the $\delta$ in (i) can be chosen without looking at $x$. It is \emph{$k$-Lipschitz} when $d(f(x),f(y))\le k\,d(x,y)$ for all $x$ and $y$. A \emph{contraction} is a Lipschitz map whose ratio $k$ is strictly less than $1$. An \emph{isometry} preserves distances.
\end{definition}

Every Lipschitz map is uniformly continuous. The map $x\mapsto x^2$ on $\R$ is continuous and not uniformly continuous, because $(n+\frac1n)^2-n^2>2$. The square root on $[0,1]$ is uniformly continuous and not Lipschitz.

\section{Completeness}

\begin{definition}
A sequence $(x_n)$ is \emph{Cauchy} when for every $\varepsilon>0$ there is an $N$ such that $d(x_n,x_m)<\varepsilon$ for all $n,m\ge N$. The space is \emph{complete} when every Cauchy sequence converges. A normed space that is complete is called a \emph{Banach space}.
\end{definition}

A convergent sequence is Cauchy. A Cauchy sequence that has a convergent subsequence converges to the same limit, because $d(x_n,x)\le d(x_n,x_{n_k})+d(x_{n_k},x)$. The line $\R$ is complete (Theorem~\ref{thm:reals}). The set $\Q$ is not, and neither is $]0,1]$: the sequence $\frac1n$ is Cauchy in $]0,1]$ and has no limit there. Figure~\ref{fig:rwopen} is the picture. Completeness is not a topological property, unlike openness. The sets $]0,1]$ and $[1,+\infty[$ are homeomorphic by $x\mapsto\frac1x$, and only the second is complete.

\begin{proposition}\label{prop:complete}
\begin{enumerate}[label=(\alph*)]
\item A closed subset of a complete space is complete. A complete subset of any metric space is closed.
\item The spaces $\R^n$ and $\C^n$ are complete for $\|\cdot\|_1$, for $\|\cdot\|_2$ and for $\|\cdot\|_\infty$.
\item The space $B(S)$ with the sup norm is a Banach space. If $S$ is a metric space, the bounded continuous functions form a closed subspace $C_b(S)$, hence a Banach space. In particular, the uniform limit of continuous functions is itself continuous.
\item In a Banach space an absolutely convergent series converges. If $\sum\|u_k\|<\infty$, then $\sum u_k$ converges and $\|\sum u_k\|\le\sum\|u_k\|$.
\end{enumerate}
\end{proposition}
\begin{proof}
(a) A Cauchy sequence in a closed set $F$ converges in $X$, and Proposition~\ref{prop:closure}(b) puts the limit in $F$. If $A$ is complete and $a_n\to x$ with $a_n\in A$, then $(a_n)$ is Cauchy, so it converges to a point of $A$. Uniqueness of limits forces $x\in A$. (b) The comparisons between the norms say that a sequence is Cauchy for one of them if and only if it is Cauchy for all of them, and likewise for convergence. For $\|\cdot\|_\infty$ this happens if and only if each coordinate is Cauchy, or convergent, in $\R$. Identify $\C^n$ with $\R^{2n}$. (c) Let $(f_n)$ be Cauchy in $B(S)$. At each point $s$ one has $|f_n(s)-f_m(s)|\le\|f_n-f_m\|_\infty$, so the numerical sequence $(f_n(s))$ is Cauchy and converges to some value $f(s)$. Given $\varepsilon$, choose $N$ with $\|f_n-f_m\|_\infty\le\varepsilon$ for $n,m\ge N$, and let $m\to\infty$. Then $|f_n(s)-f(s)|\le\varepsilon$ for every $s$ and every $n\ge N$. The limit $f$ is bounded, and the convergence $f_n\to f$ is uniform. If each $f_n$ is continuous, fix $s$ and $\varepsilon$, choose $n$ with $\|f_n-f\|_\infty<\varepsilon$, and choose $\delta$ with $|f_n(t)-f_n(s)|<\varepsilon$ whenever $d(s,t)<\delta$. Then $|f(t)-f(s)|<3\varepsilon$. Thus $C_b(S)$ is closed in $B(S)$, and complete by (a). (d) The partial sums satisfy $\|\sum_{k=m}^n u_k\|\le\sum_{k=m}^n\|u_k\|$, which is small for large $m$. They form a Cauchy sequence.
\end{proof}

\section{The completion}

\begin{lemma}[extension]\label{lem:extension}
Let $D$ be dense in $X$, let $Y$ be complete, and let $f:D\to Y$ be uniformly continuous. There is then exactly one continuous extension $\bar f:X\to Y$, and that extension is uniformly continuous. If $f$ is an isometry, so is $\bar f$.
\end{lemma}
\begin{proof}
For $x\in X$, choose $d_n\in D$ with $d_n\to x$. The sequence $(d_n)$ is Cauchy, so $(f(d_n))$ is Cauchy by uniform continuity, and it converges in $Y$. The limit does not depend on the chosen sequence. Interleave two sequences that converge to $x$: the interleaved sequence still converges to $x$, its image converges, and the two limits agree. Call the common limit $\bar f(x)$. Given $\varepsilon$, choose $\delta$ so that $d(a,b)<\delta$ implies $d(f(a),f(b))<\varepsilon$ for points of $D$. If $d(x,y)<\delta$, approximating sequences $a_n\to x$ and $b_n\to y$ eventually satisfy $d(a_n,b_n)<\delta$, so $d(\bar f(x),\bar f(y))\le\varepsilon$. The extension is uniformly continuous. Two continuous extensions that agree on $D$ agree on $X$, by Proposition~\ref{prop:closure}(c)(ii). Distances pass to the limit, which is the claim about isometries.
\end{proof}

\begin{theorem}[completion]\label{thm:completion}
Every metric space $X$ has a \emph{completion}: a complete metric space $\hat X$ together with an isometry $\iota:X\to\hat X$ whose image is dense. If $(\hat X',\iota')$ is another completion, there is a unique isometric bijection $\phi:\hat X\to\hat X'$ with $\phi\circ\iota=\iota'$.
\end{theorem}
\begin{proof}
\emph{Existence}, by Kuratowski's embedding. Assume $X\ne\varnothing$ and fix $a\in X$. For each $x\in X$, let $\iota(x)$ be the function on $X$ given by $f_x(s)=d(x,s)-d(a,s)$. The triangle inequality bounds it by $d(x,a)$, and the function is continuous, so $f_x\in C_b(X)$. Moreover $\|f_x-f_y\|_\infty=\sup_s|d(x,s)-d(y,s)|=d(x,y)$, and the supremum is attained at $s=y$. So $\iota$ is an isometry into the Banach space $C_b(X)$ of Proposition~\ref{prop:complete}(c). Let $\hat X$ be the closure of $\iota(X)$. It is complete by Proposition~\ref{prop:complete}(a), and $\iota(X)$ is dense in it. \emph{Uniqueness.} The composite $\iota'\circ\iota^{-1}:\iota(X)\to\hat X'$ is an isometry from a dense subset into a complete space. Lemma~\ref{lem:extension} extends it to an isometry $\phi:\hat X\to\hat X'$. The image of $\phi$ is complete, hence closed, and it contains the dense set $\iota'(X)$, so $\phi$ is onto. An isometry $\phi$ satisfying $\phi\circ\iota=\iota'$ is already determined by its values on the dense set $\iota(X)$. A mere map of sets that agrees with $\iota'$ on that set may do anything off it, so the uniqueness is among isometries, which is what the statement asks for.
\end{proof}

This proof cannot build $\R$, because $C_b(X)$ is made of real-valued functions. That is why Chapter~\ref{ch:sets} used Cauchy sequences. The same construction works for a general metric space: classes of Cauchy sequences of $X$, with $d([x],[y])=\lim d(x_n,y_n)$. It is a second proof of the theorem. When $X$ carries more structure, the completion inherits it. The completion of a normed space is a Banach space, because the operations are uniformly continuous on bounded sets and therefore extend. Chapter~\ref{ch:lp} finds $L^1$ and $L^2$ as completions of the continuous functions for integral norms (Exercise~\ref{ex:L1incomplete}).

\section{The \texorpdfstring{$p$}{p}-adic numbers as a completion}

In 1897 Kurt Hensel proposed an idea that looks like a joke: treat a number divisible by a high power of $p$ as small. Volume~2 names the $p$-adic absolute value in its field guide, as a preview. For the distance of Example~\ref{ex:metrics}(d) one has $p^n\to0$, the series $1+p+p^2+\cdots$ converges (to $\frac1{1-p}$), and two integers can be very close without being close on the real line. The numbers $1$ and $1+5^{10}$ differ by nearly ten million, and their $5$-adic distance is $5^{-10}$. In 1916 Alexander Ostrowski proved that, up to equivalence, the usual absolute value and the $p$-adic ones are the only absolute values on $\Q$. That theorem is quoted here and not used. Completing $\Q$ therefore produces $\R$ and the fields $\Q_p$, and nothing else. Number theory has one real place and one place for each prime, a symmetry that Volume~13 exploits. This section builds $\Q_p$ as a space, together with its ring of integers $\Z_p$. The arithmetic, Hensel's lemma in general, the squares, the extensions, belongs to Volume~13.

\begin{theorem}\label{thm:Qp}
Let $\Q_p$ be the completion of $(\Q,|\cdot|_p)$. Addition and multiplication extend from $\Q$, and $\Q_p$ is a field. The absolute value $|\cdot|_p$ extends to an ultrametric absolute value on $\Q_p$, with values in $\{p^k:k\in\Z\}\cup\{0\}$. The closed unit ball $\Z_p=\{x:|x|_p\le1\}$ is a subring. It is the closure of $\Z$, and $x\mapsto(x\bmod p^n)_{n\ge1}$ is a ring isomorphism from $\Z_p$ onto the ring of \emph{compatible sequences}: sequences $(c_n)_{n\ge1}$ with $c_n\in\Z/p^n\Z$ and $c_{n+1}\equiv c_n\pmod{p^n}$, added and multiplied term by term. Each $x\in\Z_p$ can be written in exactly one way as a convergent series $\sum_{k\ge0}a_kp^k$ with digits $a_k\in\{0,\dots,p-1\}$.
\end{theorem}
\begin{proof}
In this argument $|\cdot|$ stands for $|\cdot|_p$. Use the Cauchy-sequence model of the completion: elements of $\Q_p$ are classes of Cauchy sequences of rationals, and $\Q$ sits inside as the constant sequences. Parts (a) and (b) of Lemma~\ref{lem:cauchybasic} in Chapter~\ref{ch:sets}, and the bound $|x_n|\ge\delta$ for large $n$ in part (c), hold word for word with $|\cdot|$ in place of the usual absolute value. Their proofs used only the triangle inequality and multiplicativity. The classes therefore form a field, by the same argument as Theorem~\ref{thm:reals}. If $(x_n)$ is Cauchy and not null, then $|x_n|$ is eventually constant. Indeed $|x_n|\ge\delta$ for large $n$, and once $|x_n-x_m|<\delta$ the ultrametric inequality gives $|x_n|=|x_m|$: if $|a|>|b|$ then $|a+b|=|a|$, because $|a|\le\max(|a+b|,|b|)$. So $|x|=\lim|x_n|$ is well defined. It is multiplicative and ultrametric, by passage to the limit, its values lie in $p^\Z\cup\{0\}$, and it induces the metric of the completion.

The set $\Z_p$ is a subring, by the ultrametric inequality and multiplicativity. It is closed, being a closed ball, and it contains $\Z$. If $|x|\le1$, pick $q\in\Q$ with $|x-q|<p^{-n}$. Then $|q|\le1$, so $q=\frac ab$ with $p\nmid b$. Choose $c\in\Z$ with $bc\equiv1\pmod{p^n}$. Then $|q-ac|=|a|\,|1-bc|/|b|\le p^{-n}$. So $x$ lies within $p^{-n}$ of the integer $ac$, and $\Z$ is dense in $\Z_p$.

For $x\in\Z_p$ and $n\ge1$, let $x\bmod p^n$ be the class modulo $p^n$ of any integer $m$ with $|x-m|\le p^{-n}$. Two such integers differ by a multiple of $p^n$, so the class is well defined, and it is compatible with reduction from $p^{n+1}$ to $p^n$. The map is a ring morphism: approximate $x$ and $y$ by integers $m$ and $m'$, and then $m+m'$ and $mm'$ approximate $x+y$ and $xy$ to within $p^{-n}$. It is injective: if every reduction vanishes, then $|x|\le p^{-n}$ for every $n$. It is surjective: given a compatible sequence $(c_n)$, choose integers $m_n$ in the classes $c_n$. Then $|m_{n+1}-m_n|\le p^{-n}$, so $(m_n)$ is Cauchy, and its limit reduces to $c_n$ for each $n$.

For the digits, let $a_0\in\{0,\dots,p-1\}$ represent $x\bmod p$. Then $|x-a_0|\le p^{-1}$, so $x-a_0=py$ with $|y|\le1$, and the construction continues with $y$. Induction gives $|x-\sum_{k<n}a_kp^k|\le p^{-n}$, so the series converges to $x$. For uniqueness, if two digit sequences first differ at the index $j$, the difference of the sums is $p^j$ times $(a_j-a'_j)+p(\cdots)$, and that difference has absolute value exactly $p^{-j}$.
\end{proof}

The field guide of Volume~2 names the symbol $\Z_p$ and the absolute value $|x|_p$, and nothing more. The closed unit ball, and its identification with the compatible residue sequences, are built in this chapter. A $p$-adic integer can therefore be seen in two ways: as a compatible list of residues modulo $p$, $p^2$, $p^3$, and so on, which is pure congruence arithmetic, or as a $p$-adic limit of ordinary integers. For example $\frac1{1-p}=\sum_{k\ge0}p^k$ in $\Z_p$, and $-1=\sum_{k\ge0}(p-1)p^k$. In $\Z_2$ one has $-1=\ldots1111$, the two's complement of computers.

The digits give $\Z_p$ a shape. Sort the $p$-adic integers by the first digit. That cuts $\Z_p$ into $p$ balls of radius $\frac1p$. Sorting each of those balls by the second digit produces $p^2$ balls of radius $\frac1{p^2}$, and the process continues, like the branches of a tree. A sequence in $\Z_p$ must return infinitely often to one of the $p$ first branches, then to one of its $p$ sub-branches, and so on down the tree. The argument just given is the whole proof of the next statement. It is the Bolzano--Weierstrass theorem of $\Z_p$, with ``cut the interval in two'' replaced by ``read the next digit''.

\begin{figure}[ht]
\centering
\fig{ch02-rw-ultra}
\caption{In an ultrametric, a ball contained in a larger ball has every one of its points as a centre, and two balls are nested or disjoint. The balls of $\Z_p$ are of this kind: the ball of radius $p^{-n}$ is also the ball of radius $p^{-n+1}$ in the open sense.}
\label{fig:rwultra}
\end{figure}

\begin{corollary}[$\Z_p$ is compact]\label{cor:Zpcompact}
\begin{enumerate}[label=(\alph*)]
\item For each $n\ge1$, the ring $\Z_p$ is the disjoint union of the $p^n$ balls $a+p^n\Z_p=\{x:|x-a|_p\le p^{-n}\}$, with $a=0,1,\dots,p^n-1$. Each of those balls is both open and closed in $\Q_p$.
\item Every sequence in $\Z_p$ has a subsequence that converges to a point of $\Z_p$.
\item The same holds in every ball $x+p^k\Z_p$ of $\Q_p$, for $x\in\Q_p$ and $k\in\Z$. Thus every point of $\Q_p$ has a neighbourhood with the property that every sequence in it has a convergent subsequence. The field $\Q_p$ itself does not have the property.
\end{enumerate}
\end{corollary}
\begin{proof}
Here too, $|\cdot|$ is shorthand for $|\cdot|_p$. (a) By Theorem~\ref{thm:Qp}, each $x\in\Z_p$ has a residue $x\bmod p^n$, the class of exactly one $a\in\{0,\dots,p^n-1\}$, and that residue is the class of $a$ exactly when $|x-a|\le p^{-n}$. The $p^n$ balls therefore cover $\Z_p$ and are pairwise disjoint. Each lies in $\Z_p$ by the ultrametric inequality, since $|a|\le1$. The absolute value takes values in $p^{\Z}\cup\{0\}$, so $|x-a|\le p^{-n}$ says the same thing as $|x-a|<p^{-n+1}$. The ball is a closed ball and an open ball at once. Figure~\ref{fig:rwultra} is the same nesting in a picture.

(b) Let $(x_j)$ be a sequence in $\Z_p$. One of the $p$ balls of radius $p^{-1}$ in (a) contains $x_j$ for infinitely many $j$. Call it $B_1$. Apply (a) with $n=2$. A ball of radius $p^{-2}$ that meets $B_1$ lies inside $B_1$, because two points of such a ball are within $p^{-2}$ of each other and therefore have the same residue modulo $p$. There are finitely many of them, and one of them, $B_2$, contains $x_j$ for infinitely many indices. Continue. The result is a nested sequence $B_1\supseteq B_2\supseteq\cdots$, where $B_n$ has radius $p^{-n}$ and contains infinitely many terms. Choose $j_1<j_2<\cdots$ with $x_{j_n}\in B_n$. Whenever $m\ge n$, the two terms $x_{j_m}$ and $x_{j_n}$ both lie in $B_n$, and the ultrametric inequality gives $|x_{j_m}-x_{j_n}|\le p^{-n}$. The subsequence is Cauchy. The ring $\Z_p$ is a closed subset of the complete space $\Q_p$, hence complete by Proposition~\ref{prop:complete}(a), so the subsequence converges in $\Z_p$.

(c) The map $y\mapsto x+p^ky$ is a bijection from $\Z_p$ onto $x+p^k\Z_p=\{z:|z-x|\le p^{-k}\}$, and it multiplies every distance by $p^{-k}$. Sequences are carried to sequences, and a convergent subsequence is carried to a convergent subsequence. The ball just constructed contains the open ball $B(x,p^{-k})$, and is therefore a neighbourhood of $x$. Finally $|p^{-n}|=p^n\to\infty$. A convergent sequence is bounded, so no subsequence of $(p^{-n})$ converges.
\end{proof}

The names of these properties are given in Chapter~\ref{ch:compact}. A metric space in which every sequence has a convergent subsequence is \emph{compact} (Theorem~\ref{thm:compactmetric}), so $\Z_p$ is compact, and $\Q_p$ is \emph{locally compact}: each point has a compact neighbourhood. The line $\R$ behaves the same way, with $[-1,1]$ in the role of $\Z_p$. The resemblance stops there. The interval $[-1,1]$ is connected, and it is not a subring of $\R$. The ring $\Z_p$ falls, by (a), into $p^n$ open pieces for every $n$. Its picture is not an interval. Exercise~\ref{ex:Z2cantor} shows that the picture is the Cantor set. Compactness of $\Z_p$ is what Volume~8 needs in order to put a Haar measure on $\Q_p$.

\section{Banach's fixed point theorem}

In the summer of 1916 Hugo Steinhaus was walking in the Planty park in Krak\'ow when he overheard the words ``Lebesgue integral'' from a bench. He stopped. Two young men were talking mathematics. One of them was Stefan Banach, who had no degree and was living on odd jobs. Steinhaus later called Banach his greatest discovery. Banach's doctoral thesis, written in 1920 and published in 1922, founded the theory of complete normed spaces, now called Banach spaces, and it contained the following theorem. It may be the existence theorem that mathematics uses most often.

\begin{theorem}[Banach, 1922]\label{thm:banach}
Let $X$ be a non-empty complete metric space, and let $f:X\to X$ be a $k$-contraction with $k<1$. The map $f$ then has exactly one fixed point, written $x^*$. For every $x_0\in X$ the iterates $x_{n+1}=f(x_n)$ converge to $x^*$, and
\[
d(x_n,x^*)\le\frac{k^n}{1-k}\,d(x_0,x_1)\qquad\text{and}\qquad d(x_n,x^*)\le\frac k{1-k}\,d(x_{n-1},x_n).
\]
\end{theorem}
\begin{proof}
Induction gives $d(x_j,x_{j+1})\le k^jd(x_0,x_1)$. For $m>n$,
\[
d(x_n,x_m)\le\sum_{j=n}^{m-1}k^jd(x_0,x_1)\le\frac{k^n}{1-k}d(x_0,x_1).
\]
The sequence is Cauchy, so it converges to some $x^*$. Continuity of $f$ gives $x^*=\lim x_{n+1}=\lim f(x_n)=f(x^*)$. Letting $m\to\infty$ produces the first bound. The same bound, started at $x_{n-1}$, produces the second. If $x^*$ and $y^*$ were both fixed, then $d(x^*,y^*)=d(f(x^*),f(y^*))\le k\,d(x^*,y^*)$, which forces them to be equal.
\end{proof}

\begin{think}
Does the theorem need completeness? Does it need $k<1$, or is $d(f(x),f(y))<d(x,y)$ for $x\ne y$ enough? Try $f(x)=x+\frac1x$ on $[1,+\infty[$.
\end{think}

\begin{example}\label{ex:cosfixed}
The cosine maps $[0,1]$ into $[\cos1,1]\subseteq[0,1]$, and $|\cos x-\cos y|\le\sin1\,|x-y|$ there, by the mean value theorem (Volume~1, proved again in Chapter~\ref{ch:compact}), with $\sin1<0.85$. So $x=\cos x$ has a unique solution in $[0,1]$, the \emph{Dottie number} $0.7390851332\ldots$. The a posteriori bound says when to stop iterating. Convergence is linear, with ratio about $|\sin x^*|\approx0.67$: one new digit every six steps. Newton's method, $x\mapsto x-\frac{x-\cos x}{1+\sin x}$, is a contraction whose ratio tends to $0$ near $x^*$, and it doubles the number of digits at each step (Exercise~\ref{ex:newton}). Figure~\ref{fig:cobweb} draws both iterations.
\end{example}

\begin{figure}[ht]
\centering
\fig{ch02-cobweb}
\caption{The iteration of $x\mapsto\cos x$ from $x_0=0.1$. The cobweb spirals into the fixed point, where the graph meets the diagonal. The distances to $x^*$ shrink by a factor close to $|\cos'(x^*)|=\sin x^*\approx0.67$ at each step.}
\label{fig:cobweb}
\end{figure}

\section{Exponential, cosine, sine and \texorpdfstring{$\pi$}{pi}}

A debt has been sitting since Volume~1. The symbols $e^{i\theta}$, $\cos$, $\sin$ and $\pi$ have been used as if their properties were obvious. They are not. The school definition, ``$\pi$ is the ratio of a circle's circumference to its diameter'', assumes a length of a curve that has never been defined. Here is a cleaner road. Define $\exp$ by its power series. Let $\cos$ and $\sin$ be the real and imaginary parts of $\exp$ on the imaginary axis. \emph{Define} $\pi$ to be twice the first positive zero of $\cos$. Periodicity, the identity $\cos^2+\sin^2=1$, and even the fact that the unit circle has length $2\pi$, all follow.

The basic properties of $\exp$, $\ln$, $\cos$, $\sin$ and $\pi$ were admitted in Volume~1 and left unproved. They are defined and proved here, inside the Banach space $\C$.

\begin{lemma}[Cauchy product]\label{lem:cauchyproduct}
If $\sum a_n$ and $\sum b_n$ converge absolutely in $\C$, then $c_n=\sum_{k=0}^na_kb_{n-k}$ defines an absolutely convergent series, and $\sum c_n=\bigl(\sum a_n\bigr)\bigl(\sum b_n\bigr)$.
\end{lemma}
\begin{proof}
The partial sums of the absolute series satisfy $\sum_{n\le N}|c_n|\le\sum_{k,l\le N}|a_k||b_l|\le\sum|a_k|\sum|b_l|$, so $\sum c_n$ converges absolutely. The difference between the product of the partial sums and $\sum_{n\le N}c_n$ is the sum of $a_kb_l$ over $k,l\le N$ with $k+l>N$. In that range $k$ or $l$ exceeds $\frac N2$, so the modulus is at most $\sum_{k>N/2}|a_k|\sum_l|b_l|+\sum_k|a_k|\sum_{l>N/2}|b_l|$, which tends to $0$.
\end{proof}

\begin{definition}
For $z\in\C$, set $\exp z=\sum_{n\ge0}\frac{z^n}{n!}$. The series converges absolutely: compare it with $\sum\frac{|z|^n}{n!}$, whose terms are eventually smaller than $2^{-n}$. For real $t$, set $\cos t=\operatorname{Re}\exp(it)$ and $\sin t=\operatorname{Im}\exp(it)$, that is $\cos t=\sum\frac{(-1)^kt^{2k}}{(2k)!}$ and $\sin t=\sum\frac{(-1)^kt^{2k+1}}{(2k+1)!}$. Write $e^z=\exp z$ and $e=\exp1$.
\end{definition}

\begin{theorem}\label{thm:exp}
\begin{enumerate}[label=(\alph*)]
\item $\exp(z+w)=\exp z\exp w$, the value $\exp z$ is never $0$, $\overline{\exp z}=\exp\bar z$, and $|\exp(it)|=1$ for real $t$, so $\cos^2t+\sin^2t=1$.
\item The map $\exp$ is differentiable on $\C$, with derivative $\exp$ itself. Hence $\cos'=-\sin$ and $\sin'=\cos$ on $\R$.
\item On $\R$, the exponential is a strictly increasing bijection onto $]0,+\infty[$. The inverse $\ln$ turns products into sums, $\ln(xy)=\ln x+\ln y$, and its derivative is $\ln'(x)=\frac1x$.
\item There is a smallest $t>0$ with $\cos t=0$. Call it $\frac\pi2$. One then has $e^{i\pi/2}=i$, $e^{i\pi}=-1$ and $e^{2i\pi}=1$, so $\exp$ is periodic of period $2i\pi$.
\item For real $t$, one has $e^{it}=1$ if and only if $t\in2\pi\Z$, and $t\mapsto e^{it}$ maps $[0,2\pi[$ bijectively onto the unit circle.
\item The curve $t\mapsto e^{it}$, for $0\le t\le\theta$, has length $\theta$, once length means the integral built in Theorem~\ref{thm:riemanncont}. The length clause is postponed until Chapter~\ref{ch:integral}. Parts (a)--(e) do not use it. Thus $\pi$ is half the perimeter of the unit circle, which is the school description, now with a definition behind it.
\end{enumerate}
\end{theorem}
\begin{proof}
(a) Lemma~\ref{lem:cauchyproduct} and the binomial formula give $\exp z\exp w=\sum_n\sum_{k\le n}\frac{z^kw^{n-k}}{k!(n-k)!}=\sum_n\frac{(z+w)^n}{n!}$. In particular $\exp z\exp(-z)=\exp0=1$. Conjugation is continuous and commutes with the partial sums. Then $|e^{it}|^2=e^{it}\overline{e^{it}}=e^{it}e^{-it}=1$.

(b) Part (a) rearranges the difference quotient:
\[
\begin{aligned}\exp(z+h)-\exp z-h\exp z&=\exp z\,(\exp h-1-h),\\|\exp h-1-h|&\le\sum_{n\ge2}\frac{|h|^n}{n!}\le|h|^2e^{|h|},\end{aligned}
\]
which is $o(|h|)$. For real $t$, the chain rule for a complex function of a real variable gives $\frac{\dd}{\dd t}e^{it}=ie^{it}$: apply the estimate with $h=is$. Real and imaginary parts are the derivatives of cosine and sine.

(c) For $x>0$ one has $\exp x>1+x>1$. For $x<y$, the identity $\exp y=\exp x\exp(y-x)$ and the positivity $\exp x=(\exp\frac x2)^2>0$ give $\exp y>\exp x$. Since $\exp x\ge1+x$ and $\exp(-x)=\frac1{\exp x}$, the exponential tends to $+\infty$ at $+\infty$ and to $0$ at $-\infty$. It is continuous and strictly increasing, so Lemma~\ref{lem:monoivt} says it takes every value in $]0,+\infty[$. The inverse $\ln$ turns products into sums because $\exp$ turns sums into products. Differentiability of an inverse, from Volume~1, gives $\ln'(x)=\frac1{\exp'(\ln x)}=\frac1x$.

(d) The series gives $\sin t\ge t-\frac{t^3}6>0$ on $]0,2]$. The terms after $t$ decrease in modulus and alternate in sign, because $\frac{t^2}{(2k+2)(2k+3)}<1$. So cosine is strictly decreasing on $[0,2]$. Also $\cos2\le-\frac13$. The cosine terms do not decrease from the first term: at $t=2$ the term $\frac{t^2}2$ equals $2$, twice the constant term. From the term $\frac{t^4}{24}$ onward they do, because the ratio of successive absolute values is $\frac{t^2}{(2k+1)(2k+2)}\le\frac4{5\cdot6}<1$ for $k\ge2$. From that term on, the tail after $1-\frac{t^2}2+\frac{t^4}{24}$ is negative, and $\cos2\le1-2+\frac{2^4}{24}=-\frac13<0$. Lemma~\ref{lem:monoivt} says cosine vanishes at exactly one point of $]0,2[$, and $\cos>0$ on $[0,\frac\pi2[$. Then $\sin\frac\pi2=1$, because the sine is positive and $\sin^2=1-\cos^2$. Hence $e^{i\pi/2}=i$, and powers give $e^{i\pi}=-1$ and $e^{2i\pi}=1$. Part (a) then yields $\exp(z+2i\pi)=\exp z$.

(e) Let $0<t<2\pi$ and write $e^{it/4}=u+iv$. Since $0<\frac t4<\frac\pi2$, one has $u=\cos\frac t4>0$ and $v=\sin\frac t4>0$: sine increases on $[0,\frac\pi2]$ from $0$, because its derivative $\cos$ is positive there. Now $e^{it}=(u+iv)^4=u^4-6u^2v^2+v^4+4iuv(u^2-v^2)$. If this equalled $1$, then $u^2=v^2=\frac12$, and the real part would be $\frac14-\frac64+\frac14=-1\ne1$. So $e^{it}\ne1$ on $]0,2\pi[$. Periodicity upgrades this to: $e^{it}=1$ if and only if $t\in2\pi\Z$. Injectivity on $[0,2\pi[$ follows, because $e^{is}=e^{it}$ means $e^{i(s-t)}=1$. For surjectivity, take $|u+iv|=1$. If $u\ge0$ and $v\ge0$, cosine decreases continuously from $1$ to $0$ on $[0,\frac\pi2]$, so $u=\cos t$ for some $t$ in that interval, by Lemma~\ref{lem:monoivt}, and then $v=\sin t$, because both are non-negative and have the same square. The remaining quadrants are the first one multiplied by $i$, by $-1$, or by $-i$, which adds $\frac\pi2$, $\pi$, or $\frac{3\pi}2$.

(f) The length of a $C^1$ curve $\gamma:[a,b]\to\C$ is $\int_a^b|\gamma'(t)|\dd t$, the speed integrated over time. It is also the supremum of the lengths of inscribed polygons, a fact not needed here. This is the integral that Volume~1 admitted and did not construct. Theorem~\ref{thm:riemanncont} builds it, and parts (a)--(e) do not use it. Here $\bigl|\frac{\dd}{\dd t}e^{it}\bigr|=|ie^{it}|=1$, so the length is $\int_0^\theta 1\dd t=\theta$.
\end{proof}

The argument used one lemma, a first case of the intermediate value theorem. The general theorem is in Chapter~\ref{ch:topology}.

\begin{lemma}\label{lem:monoivt}
Let $f:[a,b]\to\R$ be continuous and strictly monotone. Then every value strictly between $f(a)$ and $f(b)$ is attained. The same holds on an open interval $]a,b[$, with $-\infty\le a<b\le+\infty$, for the values strictly between the limits of $f$ at the endpoints.
\end{lemma}
\begin{proof}
Suppose $f$ increases, and $f(a)<y<f(b)$. Let $s=\sup\{x\in[a,b]:f(x)<y\}$. If $f(s)<y$, then $s<b$, and continuity gives $f<y$ on some interval $[s,s+\eta]$, which contradicts the choice of $s$. If $f(s)>y$, then $s>a$, and $f>y$ on some $[s-\eta,s]$. Since $f$ increases, $f>y$ on $[s-\eta,b]$, so $s-\eta$ is a smaller upper bound. Hence $f(s)=y$. On an open interval, the definition of the limits supplies $a'<b'$ inside the interval with $f(a')<y<f(b')$.
\end{proof}

\section{Fractals as fixed points}

Here is a game. Draw any shape in the plane: a cat, a blob, a single point. Shrink the drawing by a factor $2$ towards each corner of a triangle, keep the three small copies, and do it again. After ten rounds the cat has gone, and the Sierpi\'nski triangle stands in its place, the same triangle no matter what was drawn first. It is Dottie's calculator again, played with shapes instead of numbers. One needs a distance between shapes. The right one is due to Hausdorff.

\begin{definition}
Let $\mathcal F(X)$ be the set of non-empty closed bounded subsets of a metric space $X$. The \emph{Hausdorff distance} between $A,B\in\mathcal F(X)$ is
\[
d_H(A,B)=\max\Bigl(\sup_{a\in A}d(a,B),\ \sup_{b\in B}d(b,A)\Bigr),
\]
the smallest $r$ such that each set lies in the closed $r$-neighbourhood of the other.
\end{definition}

This is a metric on $\mathcal F(X)$. The equality $d_H(A,B)=0$ forces $A\subseteq\bar B=B$ and $B\subseteq A$. The triangle inequality follows from $d(a,C)\le d(a,b)+d(b,C)$: take the infimum over $b\in B$, then the suprema.

\begin{theorem}\label{thm:hausdorffcomplete}
Whenever $X$ is complete, the space $(\mathcal F(X),d_H)$ is complete as well.
\end{theorem}
\begin{proof}
Let $(A_n)$ be Cauchy, and set $A=\bigcap_n\overline{\bigcup_{k\ge n}A_k}$, a closed set. Fix $\varepsilon>0$ and choose $N_0<N_1<\cdots$ with $d_H(A_m,A_n)<\varepsilon2^{-j-1}$ whenever $m,n\ge N_j$. \emph{Points of $A_n$ are near $A$.} Take $n\ge N_0$ and $x\in A_n$. Set $x_0=x$ and $m_0=n$. For $j\ge1$ let $m_j=\max(n,N_j)$ and choose $x_j\in A_{m_j}$ with $d(x_{j-1},x_j)<\varepsilon2^{-j}$. This is possible because $m_{j-1}$ and $m_j$ are at least $N_{j-1}$. The sequence $(x_j)$ is Cauchy and converges to some $y$ with $d(x,y)<\varepsilon$, and $y\in A$, because for each $p$ it is a limit of points $x_j\in A_{m_j}$ with $m_j\ge p$ eventually. So $A$ is non-empty and $\sup_{x\in A_n}d(x,A)\le\varepsilon$. \emph{Points of $A$ are near $A_n$.} Take $n\ge N_0$ and $a\in A$. Since $a\in\overline{\bigcup_{k\ge n}A_k}$, there exist $k\ge n$ and $y\in A_k$ with $d(a,y)<\frac\varepsilon2$, and $d(y,A_n)\le d_H(A_k,A_n)<\frac\varepsilon2$. So $d(a,A_n)<\varepsilon$. Finally $A$ is bounded, because it lies within $\varepsilon$ of the bounded set $A_n$. Hence $A\in\mathcal F(X)$ and $d_H(A_n,A)\le\varepsilon$ for $n\ge N_0$.
\end{proof}

\begin{theorem}[Hutchinson, 1981]\label{thm:hutchinson}
Let $X$ be complete, and let $f_1,\dots,f_m:X\to X$ be contractions with ratios $k_i<1$. There is a unique $K\in\mathcal F(X)$ with $K=\overline{f_1(K)\cup\cdots\cup f_m(K)}$, the \emph{attractor} of the system. For every $A\in\mathcal F(X)$ the iterates $F^n(A)$ of $F(A)=\overline{\bigcup_if_i(A)}$ converge to $K$ in the Hausdorff distance.
\end{theorem}
\begin{proof}
The set $F(A)$ is closed, non-empty and bounded, so $F$ maps $\mathcal F(X)$ into itself. For each index $i$ one has $d(f_i(a),f_i(B))\le k_id(a,B)$, and therefore $d_H(f_i(A),f_i(B))\le k_id_H(A,B)$, with $d_H$ extended to bounded sets by the same formula. A point of $\bigcup f_i(A)$ lies in some $f_i(A)$, within $k_id_H(A,B)$ of $f_i(B)$. Therefore $d_H\bigl(\bigcup f_i(A),\bigcup f_i(B)\bigr)\le\max_ik_i\,d_H(A,B)$. Closures do not change $d_H$, because $d(x,\bar B)=d(x,B)$ and $d(\cdot,B)$ is continuous. The map $F$ is therefore a contraction of the complete space $\mathcal F(X)$, and Theorem~\ref{thm:banach} applies.
\end{proof}

For $f_1(x)=\frac x3$ and $f_2(x)=\frac{x+2}3$ on $\R$, the attractor is the \emph{Cantor set} $C$, the numbers in $[0,1]$ whose ternary expansion uses only the digits $0$ and $2$. The map $f_1$ prepends a digit $0$ and $f_2$ prepends a digit $2$, so $f_1(C)\cup f_2(C)=C$, and $C$ is closed and bounded (Exercise~\ref{ex:cantorifs}). Starting from $[0,1]$, the iterate $F^n([0,1])$ is the union of the $2^n$ intervals of stage $n$. Starting from a single point, it is a set of $2^n$ points, and that set converges to the same $C$. The same phenomenon, drawn in the plane, is Figure~\ref{fig:sierpinski}.

\begin{figure}[ht]
\centering
\fig{ch02-sierpinski}
\caption{Three contractions of ratio $\frac12$, towards the corners of a triangle. Starting from a square, the iterates $F(A)$, $F^2(A)$ and $F^6(A)$ forget the square and converge to the Sierpi\'nski triangle, the attractor of Theorem~\ref{thm:hutchinson}.}
\label{fig:sierpinski}
\end{figure}

\begin{beauty}
Banach's theorem has a proof of a few lines, and it solves equations of every kind: a numerical equation such as $x=\cos x$, an algebraic equation such as Newton's method and Hensel's lemma in $\Z_p$ (Exercise~\ref{ex:hensel}), a differential equation (Picard--Lindel\"of, Exercise~\ref{ex:picard}), a probabilistic equation (PageRank, and the stationary laws of Markov chains, as numerical schemes), and a geometric equation (a fractal). Each time the work is to find a complete space on which the problem is a contraction. Completeness itself, the property that made $\R$ out of $\Q$, makes $\Q_p$ out of $\Q$ for another distance. Ostrowski's theorem, quoted above, says that $\R$ and the fields $\Q_p$ are the only completions of $\Q$ for an absolute value.
\end{beauty}

\section{What we proved, and what comes next}

The calculator stops at $0.739085\ldots$ because cosine is a contraction of the complete interval $[0,1]$, and a contraction of a complete space has exactly one fixed point, reached from anywhere (Theorem~\ref{thm:banach}). The Cantor set appears for the same reason, one level up. The two maps act on the complete space of closed bounded sets, and there they form a contraction (Theorems~\ref{thm:hausdorffcomplete} and~\ref{thm:hutchinson}).

Completeness carried the chapter. Theorem~\ref{thm:completion} says that every metric space has a completion, and that the completion is unique up to isometry. The completion of $\Q$ for $|\cdot|_p$ is the field $\Q_p$. Its unit ball $\Z_p$ is a ring, every element has a unique expansion in digits, and the ball is compact (Theorem~\ref{thm:Qp}, Corollary~\ref{cor:Zpcompact}). One debt to Volume~1 is paid: $\exp$, $\cos$, $\sin$ and $\pi$ are defined, and their properties are proved (Theorem~\ref{thm:exp}).

Proposition~\ref{prop:closure}(c) left a question. Continuity can be tested on open sets alone, with no mention of the distance. How much of this chapter survives if the distance is forgotten and only the open sets are kept? Chapter~\ref{ch:topology} starts from that question.

\section{Exercises}

\begin{exercise}[$\star$]\label{ex:ultra}
Let $d$ be an ultrametric, so $d(x,z)\le\max(d(x,y),d(y,z))$. Show that every triangle is isosceles, with the two longest sides equal. Show that every point of a ball is a centre of that ball, that two balls are either disjoint or nested, and that every open ball is also closed.
\end{exercise}

\begin{exercise}[$\star$]\label{ex:newton}
Let $a>0$ and $N(x)=\frac12\bigl(x+\frac ax\bigr)$. Show that $N$ maps $[\sqrt a,+\infty[$ into itself and is a contraction there with ratio $\frac12$, and that the errors $e_n=x_n-\sqrt a$ satisfy $e_{n+1}=\frac{e_n^2}{2x_n}$. Starting from $x_0=2$ with $a=2$, how many correct digits does $x_5$ have?
\end{exercise}

\begin{exercise}[$\star$]\label{ex:cantorifs}
Show that the Cantor set $C$ is closed and bounded and satisfies $C=f_1(C)\cup f_2(C)$ for $f_1(x)=\frac x3$ and $f_2(x)=\frac{x+2}3$. Name the attractor of the two maps $\frac x2$ and $\frac{x+1}2$. Of $\frac x3$, $\frac{x+1}3$ and $\frac{x+2}3$?
\end{exercise}

\begin{exercise}[$\star\star$]\label{ex:L1incomplete}
On $C([0,1])$, set $\|f\|_1=\int_0^1|f|$. Show that this is a norm, and that $C([0,1])$ is not complete for it. Find a Cauchy sequence of continuous functions that converges, in a sense you must make precise, to the indicator of $[\frac12,1]$.
\end{exercise}

\begin{exercise}[$\star\star$]\label{ex:cantorinter}
Show that $X$ is complete if and only if every decreasing sequence $F_1\supseteq F_2\supseteq\cdots$ of non-empty closed sets with $\operatorname{diam}F_n\to0$ has non-empty intersection (Cantor's intersection theorem). Prove that the intersection is a single point, and give an example showing that the hypothesis on the diameters cannot be omitted.
\end{exercise}

\begin{exercise}[$\star\star$]\label{ex:picard}
\emph{Picard--Lindel\"of.} Let $f:[0,T]\times\R\to\R$ be continuous and $L$-Lipschitz in its second variable, and let $y_0\in\R$. On $C([0,T])$, consider $\Phi(y)(t)=y_0+\int_0^tf(s,y(s))\dd s$ and the norm $\|y\|_L=\sup_te^{-2Lt}|y(t)|$. Show that $\Phi$ is a $\frac12$-contraction for $\|\cdot\|_L$, and deduce that $y'=f(t,y)$, $y(0)=y_0$, has a unique solution on $[0,T]$. For the equation $y'=y$ with $y(0)=1$, what are the iterates?
\end{exercise}

\begin{exercise}[$\star\star$]\label{ex:hensel}
In $\Z_5$, let $T(x)=x-4(x^2+1)$. Show that $T$ maps $D=\{x:|x-2|_5\le\frac15\}$ into itself and is a contraction there with ratio $\frac15$. Deduce that $-1$ has a square root in $\Z_5$ congruent to $2$ modulo $5$. Find the first three digits in its $5$-adic expansion.
\end{exercise}

\begin{exercise}[$\star\star\star$]\label{ex:devil}
\emph{The devil's staircase as a fixed point.} On the Banach space $B([0,1])$, define $\Phi(f)(x)=\frac12f(3x)$ for $x\in[0,\frac13]$, $\Phi(f)(x)=\frac12$ for $x\in\,]\frac13,\frac23[$, and $\Phi(f)(x)=\frac12+\frac12f(3x-2)$ for $x\in[\frac23,1]$. Show that $\Phi$ is a $\frac12$-contraction, that its fixed point $c$ is continuous and non-decreasing with $c(0)=0$ and $c(1)=1$, and that $c$ is constant on each interval removed from $[0,1]$ in the construction of the Cantor set.
\end{exercise}

\begin{exercise}[$\star\star$]\label{ex:Z2cantor}
\emph{A picture of $\Z_2$.} For $x=\sum_{k\ge0}a_k2^k\in\Z_2$ with digits $a_k\in\{0,1\}$, let $\phi(x)=\sum_{k\ge0}2a_k3^{-k-1}$, a real number. (a) Prove that $\phi$ pairs $\Z_2$ bijectively with the Cantor set $C$. (b) Show that $\frac13|x-y|_2^{\log_23}\le|\phi(x)-\phi(y)|\le|x-y|_2^{\log_23}$ for all $x,y\in\Z_2$. (c) Conclude that convergence in $\Z_2$ is the same thing as convergence of the image in $C$. (d) Where do $-1$ and $\frac13$ land?
\end{exercise}

\begin{exercise}[$\star$]\label{ex:v7c02a}
The map $f(x)=\frac x2$ sends $(0,1]$ into itself and is a contraction. Show that it has no fixed point in $(0,1]$. Which hypothesis of Theorem~\ref{thm:banach} fails?
\end{exercise}

\begin{exercise}[$\star$]\label{ex:v7c02b}
On $\R$, let $d(x,y)=|\arctan x-\arctan y|$. Show that $d$ is a metric, that it determines the same open sets as $|x-y|$, and that $(\R,d)$ is not complete. Exhibit a $d$-Cauchy sequence with no limit in $\R$.
\end{exercise}

\begin{exercise}[$\star\star$]\label{ex:v7c02c}
Let $X$ be a metric space and set $d'(x,y)=\min(1,d(x,y))$. Show that $d'$ is a metric, that it has the same open sets as $d$, and that a sequence is $d$-Cauchy if and only if it is $d'$-Cauchy. Deduce that $X$ is complete for $d$ if and only if it is complete for $d'$.
\end{exercise}

\begin{exercise}[$\star\star$]\label{ex:v7c02d}
On $C([0,1])$, the norms $\|f\|_\infty$ and $\|f\|_1=\int_0^1|f|$ are not Lipschitz equivalent. Show that there is no constant $C$ with $\|f\|_\infty\le C\|f\|_1$ for every continuous $f$. (The other comparison, $\|f\|_1\le\|f\|_\infty$, is immediate.)
\end{exercise}

\section{Hints}

\begin{hint}{ex:ultra}
If $d(x,y)<d(y,z)$, apply the ultrametric inequality to $d(y,z)$ with the intermediate point $x$. For the last claim, a point at distance at least $r$ from the centre has a whole ball around it that stays outside.
\end{hint}
\begin{hint}{ex:newton}
The identity $N(x)-\sqrt a=\frac{(x-\sqrt a)^2}{2x}$ is the error formula, and $N'(x)=\frac12\bigl(1-\frac a{x^2}\bigr)$ lies in $[0,\frac12[$.
\end{hint}
\begin{hint}{ex:cantorifs}
Write $C=\bigcap C_n$, with $C_n$ the union of $2^n$ closed intervals. The other two attractors are intervals.
\end{hint}
\begin{hint}{ex:L1incomplete}
Take $f_n$ equal to $0$ on $[0,\frac12-\frac1n]$, equal to $1$ on $[\frac12,1]$, and affine in between. Any limit $g$ that stayed in $C([0,1])$ would have to satisfy $\int|g-\mathbf 1_{[1/2,1]}|=0$.
\end{hint}
\begin{hint}{ex:cantorinter}
For the converse, given a Cauchy sequence, use the closed tails $F_n=\overline{\{x_k:k\ge n\}}$. For a counterexample, try $F_n=[n,+\infty[$.
\end{hint}
\begin{hint}{ex:picard}
The difference satisfies $|\Phi(y)(t)-\Phi(z)(t)|\le L\int_0^te^{2Ls}\|y-z\|_L\dd s$.
\end{hint}
\begin{hint}{ex:hensel}
On $D$, write $x=2+5u$ with $|u|\le1$. Compute $T(x)-T(x')=(x-x')\bigl(1-4(x+x')\bigr)$ and look at $1-4(x+x')$ modulo $5$.
\end{hint}
\begin{hint}{ex:devil}
On each of the three pieces, $\Phi(f)-\Phi(g)$ is bounded by $\frac12\|f-g\|_\infty$. The continuous non-decreasing functions with $f(0)=0$ and $f(1)=1$ form a closed set that $\Phi$ preserves.
\end{hint}
\begin{hint}{ex:Z2cantor}
(b) If the digits of $x$ and $y$ first differ at index $j$, then $|x-y|_2=2^{-j}$, and $\phi(x)$, $\phi(y)$ lie in two different intervals of stage $j+1$ of the Cantor set, inside the same interval of stage $j$. (d) In $\Z_2$ one has $\frac13=1-\frac23=1+2\sum_{k\ge0}4^k$.
\end{hint}
\begin{hint}{ex:v7c02a}
A fixed point would solve $\frac x2=x$. The missing hypothesis is completeness of $(0,1]$.
\end{hint}
\begin{hint}{ex:v7c02b}
The sequence of integers $n=1,2,3,\ldots$ is $d$-Cauchy, because $\arctan n$ converges.
\end{hint}
\begin{hint}{ex:v7c02c}
Balls of radius less than $1$ are the same for $d$ and for $d'$. A Cauchy sequence is eventually of small steps, smaller than $1$.
\end{hint}
\begin{hint}{ex:v7c02d}
Take a tent of height $n$ and width $\frac1n$. The integral stays bounded while the height grows.
\end{hint}

\section{Solutions}

\begin{solution}{ex:ultra}
\idea One comparison does all the work: if $d(x,y)\ne d(y,z)$, then $d(x,z)$ equals the larger of the two. Once that comparison is known, balls stack the way the levels of a tree do. Figure~\ref{fig:rwultra} is the picture.
\writeup Suppose $d(x,y)<d(y,z)$. The ultrametric inequality gives $d(y,z)\le\max(d(y,x),d(x,z))$, and the maximum cannot be $d(x,y)$, so $d(y,z)\le d(x,z)\le\max(d(x,y),d(y,z))=d(y,z)$. Thus $d(x,z)=d(y,z)$. The two longest sides are equal. If $y\in B(x,r)$ and $d(z,y)<r$, then $d(z,x)\le\max(d(z,y),d(y,x))<r$, and the symmetric argument gives $B(y,r)=B(x,r)$. If two balls $B(x,r)$ and $B(x',r')$, with $r\le r'$, share a point $z$, then $B(x,r)=B(z,r)\subseteq B(z,r')=B(x',r')$. Finally, if $d(y,x)\ge r$ and $d(z,y)<r$, then $d(z,x)=d(y,x)\ge r$, because $d(z,y)<d(y,x)$. The ball $B(y,r)$ lies outside $B(x,r)$, so the open ball $B(x,r)$ is closed.
\end{solution}

\begin{solution}{ex:newton}
\idea Newton's map is a contraction, and the ratio is the derivative of $N$. At the fixed point that derivative vanishes, so the contraction becomes stronger as the iterates approach the root.
\writeup The algebra $N(x)-\sqrt a=\frac{x^2-2x\sqrt a+a}{2x}=\frac{(x-\sqrt a)^2}{2x}\ge0$ is the identity $e_{n+1}=\frac{e_n^2}{2x_n}$, and it shows $N([\sqrt a,+\infty[)\subseteq[\sqrt a,+\infty[$. For $x\ge\sqrt a$ one has $0\le N'(x)=\frac12\bigl(1-\frac a{x^2}\bigr)<\frac12$, so the mean value theorem makes $N$ a $\frac12$-contraction on this closed, hence complete, set. For $a=2$ and $x_0=2$, the error starts at $e_0\le0.586$, and $e_{n+1}\le\frac{e_n^2}{2\sqrt2}$. Rounding each bound up gives $e_1\le0.122$, $e_2\le0.0053$, $e_3\le1.0\cdot10^{-5}$, $e_4\le3.6\cdot10^{-11}$ and $e_5\le4.6\cdot10^{-22}$. So $x_5$ has at least $21$ correct digits. The true error is about $9\cdot10^{-25}$, which is $23$ matching decimals, and $e_1=\frac32-\sqrt2\approx0.086$ is smaller than the bound above. The number of digits doubles at each step. The crude ratio $\frac12$ would promise only one binary digit per step.
\end{solution}

\begin{solution}{ex:cantorifs}
\idea Write $C$ through its ternary digits. The two maps prepend a digit. For the other systems, guess the attractor by applying the maps once to $[0,1]$, and invoke uniqueness.
\writeup Let $C_0=[0,1]$ and $C_{n+1}=f_1(C_n)\cup f_2(C_n)$. Induction shows that $C_n$ is the disjoint union of the $2^n$ closed intervals $I_d=\bigl[\sum_{k\le n}d_k3^{-k},\sum_{k\le n}d_k3^{-k}+3^{-n}\bigr]$ for $d\in\{0,2\}^n$, and that $C_{n+1}\subseteq C_n$. The intersection $\bigcap_n C_n$ equals $C$. If $x=\sum_k d_k3^{-k}$ with every $d_k\in\{0,2\}$, then $x$ lies in the interval of digits $(d_1,\dots,d_n)$ at every stage. Conversely, if $x$ lies in every $C_n$, then at each stage it lies in exactly one interval, the intervals being disjoint, and the interval of stage $n+1$ sits inside the interval of stage $n$. The digits are therefore consistent, and $|x-\sum_{k\le n}d_k3^{-k}|\le3^{-n}$ gives the series. So $C$ is closed, as an intersection of closed sets, and bounded. The maps $f_i$ are injective, so $f_i(\bigcap C_n)=\bigcap f_i(C_n)$. Also $f_1(C_n)\subseteq[0,\frac13]$ and $f_2(C_n)\subseteq[\frac23,1]$ are disjoint, hence $f_1(C)\cup f_2(C)=\bigcap_n\bigl(f_1(C_n)\cup f_2(C_n)\bigr)=\bigcap_n C_{n+1}=C$. Uniqueness says $C$ is the attractor. For the maps $\frac x2$ and $\frac{x+1}2$, the interval $[0,1]=[0,\frac12]\cup[\frac12,1]$ is fixed, so it is the attractor: binary expansions use every digit. For the three maps of ratio $\frac13$, one has $[0,\frac13]\cup[\frac13,\frac23]\cup[\frac23,1]=[0,1]$ again. A fractal attractor needs missing pieces.
\end{solution}

\begin{solution}{ex:L1incomplete}
\idea Continuous functions can approximate a jump in the mean. The jump itself is not continuous, so the limit falls outside the space.
\writeup If $f\ne0$ is continuous, then $|f|>\frac{|f(x_0)|}2$ on an interval about some $x_0$, so $\int|f|>0$. Homogeneity and the triangle inequality come from those of the absolute value, together with monotonicity of the integral. For the functions $f_n$ of the hint and for $m>n$, the functions $f_n$ and $f_m$ differ only on $[\frac12-\frac1n,\frac12]$, where both take values in $[0,1]$, so $\|f_n-f_m\|_1\le\frac1n$. The sequence is Cauchy. Suppose the limit $g$ of $(f_n)$ in $\|\cdot\|_1$ were continuous. For fixed $\eta>0$, every $f_n$ equals $1$ on $[\frac12,1]$, so $\int_{1/2}^1|1-g|\le\|f_n-g\|_1\to0$, and continuity forces $g=1$ there. On $[0,\frac12-\eta]$, one has $f_n=0$ for $n>\frac1\eta$, so $g=0$ there. Since $\eta$ is arbitrary, $g=0$ on $[0,\frac12[$ and $g(\frac12)=1$, which contradicts continuity. The completion is $L^1([0,1])$ (Chapter~\ref{ch:lp}), where the indicator $\mathbf 1_{[1/2,1]}$ lives.
\end{solution}

\begin{solution}{ex:cantorinter}
\idea Shrinking closed sets pin down a Cauchy sequence. Conversely, the tails of a Cauchy sequence are shrinking closed sets.
\writeup Let $X$ be complete, and let $(F_n)$ be as in the statement. Pick $x_n\in F_n$. For $m\ge n$ one has $x_m\in F_n$, so $d(x_n,x_m)\le\operatorname{diam}F_n\to0$. The sequence $(x_n)$ is Cauchy, and it has a limit $x$. For each $n$, the point $x$ is a limit of points $x_m\in F_n$ with $m\ge n$, so $x\in F_n$. Two points of the intersection are at distance at most $\operatorname{diam}F_n$ for every $n$, so they coincide. Conversely, let $(x_k)$ be Cauchy and set $F_n=\overline{\{x_k:k\ge n\}}$. Then $\operatorname{diam}F_n=\sup_{k,l\ge n}d(x_k,x_l)\to0$, the closure not changing the diameter, so some $x$ lies in every $F_n$. For $k\ge n$ one has $d(x,x_k)\le\operatorname{diam}F_n$, and $x_k\to x$. Without a bound on the diameters, the sets $F_n=[n,+\infty[$ in $\R$ are closed, non-empty and nested, and their intersection is empty.
\end{solution}

\begin{solution}{ex:picard}
\idea The integral equation $y=\Phi(y)$ is the differential equation in another form. The weight $e^{-2Lt}$ is chosen so that the Lipschitz constant $L$ is beaten on the whole interval, not only on a short piece of it.
\writeup The space $C([0,T])$ with $\|\cdot\|_L$ is complete: $e^{-2LT}\|y\|_\infty\le\|y\|_L\le\|y\|_\infty$, so the Cauchy sequences are the same as for the sup norm. The image $\Phi(y)$ is continuous, being an integral of a continuous function, and
\[
\begin{aligned}|\Phi(y)(t)-\Phi(z)(t)|&\le\int_0^t L|y(s)-z(s)|\dd s\eqbr\le L\|y-z\|_L\int_0^t e^{2Ls}\dd s\eqbr\le\frac{e^{2Lt}}2\|y-z\|_L.\end{aligned}
\]
Multiply by $e^{-2Lt}$ and take the supremum: $\|\Phi(y)-\Phi(z)\|_L\le\frac12\|y-z\|_L$. Banach's theorem gives a unique fixed point $y$. A fixed point is of class $C^1$, with $y'=f(t,y)$ and $y(0)=y_0$, by the fundamental theorem of calculus. Conversely, any solution satisfies $y=\Phi(y)$ by integration, so the solution is unique. For $y'=y$ and $y_0=1$, the iterates that start at the constant function $1$ are $\sum_{k\le n}\frac{t^k}{k!}$. Picard's iteration rebuilds the exponential series.
\end{solution}

\begin{solution}{ex:hensel}
\idea This is Newton's method for $x^2+1$, with the derivative frozen at the first approximation. One has $2\cdot2=4$, and the inverse of $4$ modulo $5$ is $4$. Freezing the derivative loses quadratic convergence, but it keeps a contraction, which is all Banach's theorem needs.
\writeup The set $D$ is a closed ball of the complete space $\Z_5$, hence complete. For $x\in D$ one has $x^2+1\equiv4+1\equiv0\pmod5$, so $|x^2+1|_5\le\frac15$ and $|T(x)-2|_5\le\max\bigl(|x-2|_5,|4(x^2+1)|_5\bigr)\le\frac15$. Thus $T(D)\subseteq D$. Next, $T(x)-T(x')=(x-x')\bigl(1-4(x+x')\bigr)$, and $x+x'\equiv4\pmod5$, so $1-4(x+x')\equiv1-16\equiv0\pmod5$. Hence $|T(x)-T(x')|_5\le\frac15|x-x'|_5$. A fixed point satisfies $4(x^2+1)=0$, so $x^2=-1$, and $x\equiv2\pmod5$. The a priori bound gives $|T^n(2)-x|_5\le5^{-n-1}$, so $T^n(2)\equiv x\pmod{5^{n+1}}$, and an integer congruent to $T^n(2)$ modulo $5^{n+1}$ may replace it: the image moves by at most a fifth of that distance. Iterating from $2$ gives $T(2)=2-20=-18\equiv7\pmod{25}$, so $x\equiv7=2+1\cdot5\pmod{25}$. Then $T(7)=7-200=-193\equiv57\pmod{125}$, and $57=2+1\cdot5+2\cdot25$. Check: $57^2+1=3250=26\cdot125$. The digits start $2,1,2$.
\end{solution}

\begin{solution}{ex:devil}
\idea Self-similarity is a fixed-point equation. On a space of bounded functions with the sup norm, halving the values makes the map a contraction. A property that $\Phi$ preserves, and that is closed under uniform limits, passes to the fixed point.
\writeup On each of the three pieces, $|\Phi(f)(x)-\Phi(g)(x)|$ is either $0$ or $\frac12|f(y)-g(y)|$ for some $y$. So $\|\Phi(f)-\Phi(g)\|_\infty\le\frac12\|f-g\|_\infty$, and Banach's theorem in $B([0,1])$ gives a unique fixed point $c$. Let $M$ be the set of continuous non-decreasing functions with $f(0)=0$ and $f(1)=1$. It is closed in $B([0,1])$: a uniform limit keeps continuity, monotonicity, and the values at the endpoints. It is non-empty, because it contains $f(x)=x$, and $\Phi(M)\subseteq M$. On $[0,\frac13]$, the image $\Phi(f)$ rises from $0$ to $\frac12f(1)=\frac12$. It equals $\frac12$ in the middle. It rises from $\frac12+\frac12f(0)=\frac12$ to $1$ on $[\frac23,1]$. The pieces join continuously and monotonically. The iterates of any $f\in M$ stay in $M$ and converge to $c$, so $c\in M$. Finally $c=\frac12$ on $]\frac13,\frac23[$. The equation $c(x)=\frac12c(3x)$ says that $c$ is constant on each removed interval of $[0,\frac13]$ that is the image of a removed interval under $x\mapsto\frac x3$, and likewise on $[\frac23,1]$. Induction on the stage shows that $c$ is constant on every removed interval. The function $c$ is the distribution function of Volume~4, Exercise~5.7. Chapter~\ref{ch:radon} shows that its derivative vanishes almost everywhere.
\end{solution}

\begin{solution}{ex:Z2cantor}
\idea Both spaces are the set of all infinite strings of two symbols. What $\phi$ keeps is the string of digits. What it changes is the rule that turns a string into a number. Strings that share a long common prefix are close in both metrics, and the first index where they differ is the distance in both.
\writeup (a) Theorem~\ref{thm:Qp} says each $x\in\Z_2$ has exactly one digit sequence $(a_k)$, and every sequence of digits occurs: the partial sums are Cauchy, because $|2^k|_2\to0$. Exercise~\ref{ex:cantorifs} says $C$ is the set of sums $\sum_{k\ge1}d_k3^{-k}$ with $d_k\in\{0,2\}$, and the digits $d_k$ are determined by the point, which lies in exactly one interval of each stage. Therefore $\phi$, the map that turns the digits $(a_k)$ into the ternary digits $d_{k+1}=2a_k$, is a bijection. (b) Let $x\ne y$, and let $j$ be the first index where the digits differ. Then $x-y=2^j(\pm1+2z)$ with $z\in\Z_2$, so $|x-y|_2=2^{-j}$ and $|x-y|_2^{\log_23}=3^{-j}$. On the other side,
\[\phi(x)-\phi(y)=\pm2\cdot3^{-j-1}+\sum_{k>j}2(a_k-a_k')3^{-k-1},\qquad\Bigl|\sum_{k>j}2(a_k-a_k')3^{-k-1}\Bigr|\le\sum_{k>j}2\cdot3^{-k-1}=3^{-j-1}.\]
Hence $3^{-j-1}\le|\phi(x)-\phi(y)|\le3\cdot3^{-j-1}=3^{-j}$, which is the claim. (c) Part (b) says that $|x_n-x|_2\to0$ if and only if $|\phi(x_n)-\phi(x)|\to0$. So $\phi$ and $\phi^{-1}$ are continuous. As far as convergence is concerned, $\Z_2$ is the Cantor set. In particular Corollary~\ref{cor:Zpcompact}(b) says that every sequence in the Cantor set has a convergent subsequence, which Chapter~\ref{ch:compact} proves again from another side. (d) The number $-1$ has every digit equal to $1$, so $\phi(-1)=\sum_{k\ge0}2\cdot3^{-k-1}=1$, the right endpoint of $C$. For $\frac13$: in $\Z_2$ one has $\sum_{k\ge0}4^k=\frac1{1-4}=-\frac13$, so $\frac13=1-\frac23=1+2\sum_{k\ge0}4^k=1+\sum_{k\ge0}2^{2k+1}$. The digits are $a_0=1$, then $a_k=1$ for odd $k$ and $a_k=0$ for even $k\ge2$: the string $1,1,0,1,0,1,0,\ldots$. Therefore
\[\phi\bigl(\tfrac13\bigr)=\tfrac23+\sum_{m\ge0}2\cdot3^{-(2m+1)-1}=\tfrac23+\tfrac29\cdot\tfrac1{1-\frac19}=\tfrac23+\tfrac14=\tfrac{11}{12}.\]
\end{solution}

\begin{solution}{ex:v7c02a}
\idea The map shrinks distances by $\frac12$, so the contraction hypothesis holds. The space $(0,1]$ has a hole at $0$, and the iterates fall into that hole.
\writeup For $x,y\in(0,1]$ one has $\bigl|\frac x2-\frac y2\bigr|=\frac12|x-y|$, so $f$ is a $\frac12$-contraction, and $f$ sends the interval into itself. A fixed point would satisfy $\frac x2=x$, hence $x=0$, which is not in $(0,1]$. Theorem~\ref{thm:banach} does not apply, because $(0,1]$ is not complete: the iterates of $1$ are $2^{-n}$, a Cauchy sequence with no limit in $(0,1]$.
\end{solution}

\begin{solution}{ex:v7c02b}
\idea The arctangent folds the whole line into a bounded interval. Sequences that run off to infinity become Cauchy for $d$, and they have nowhere to land in $\R$ for that metric.
\writeup The difference of arctangents is symmetric, vanishes only when $x=y$ because arctangent is injective, and the triangle inequality is the ordinary one on the image. A $d$-ball of radius $\varepsilon$ about $x$ is an ordinary open interval, because arctangent is a homeomorphism from $\R$ onto $\bigl(-\frac\pi2,\frac\pi2\bigr)$. So the open sets agree. The sequence $n=1,2,3,\ldots$ satisfies $d(n,m)=|\arctan n-\arctan m|\to0$ as $n,m\to\infty$, so it is $d$-Cauchy. It does not converge in $\R$: a $d$-limit $L$ would have $\arctan n\to\arctan L$, hence $\arctan L=\frac\pi2$, which is impossible.
\end{solution}

\begin{solution}{ex:v7c02c}
\idea Cutting the metric at height $1$ changes only the large distances. Cauchy sequences, and open sets, are decided by small distances.
\writeup Symmetry and positive-definiteness pass from $d$ to $d'$. For the triangle inequality, $d'(x,z)\le 1$ and $d'(x,z)\le d(x,z)\le d(x,y)+d(y,z)$. If both $d(x,y)$ and $d(y,z)$ are at most $1$, then $d'(x,y)+d'(y,z)=d(x,y)+d(y,z)\ge d'(x,z)$. If either distance exceeds $1$, then $d'(x,y)+d'(y,z)\ge 1\ge d'(x,z)$. Balls of radius $r<1$ coincide for the two metrics, so the open sets coincide. A sequence is Cauchy when its steps are eventually smaller than every $\varepsilon$. For $\varepsilon<1$ the two metrics agree on those steps, so the Cauchy sequences are the same, and so are the convergent sequences. Completeness passes from one metric to the other.
\end{solution}

\begin{solution}{ex:v7c02d}
\idea A tall thin spike has a large supremum and a small integral. No constant can turn the integral back into the height.
\writeup Let $f_n$ be the tent that rises linearly from $0$ to $n$ on $[0,\frac1{2n}]$, falls linearly from $n$ to $0$ on $[\frac1{2n},\frac1n]$, and vanishes on $[\frac1n,1]$. Then $\|f_n\|_\infty=n$ and $\int_0^1 f_n=\frac12\cdot\frac1n\cdot n=\frac12$. If a constant $C$ satisfied $\|f\|_\infty\le C\|f\|_1$ for every continuous $f$, then $n\le C/2$ for every $n$, which is absurd. The comparison $\|f\|_1\le\|f\|_\infty$ is the integral of a function bounded by its supremum.
\end{solution}

\section*{Sources and further reading}
\addcontentsline{toc}{section}{Sources and further reading}

\begin{itemize}
\item Banach's thesis is \emph{Sur les op\'erations dans les ensembles abstraits et leur application aux \'equations int\'egrales}, Fundamenta Mathematicae 3 (1922), 133--181. It introduces Banach spaces and the fixed point theorem.
\item Fr\'echet's thesis is \emph{Sur quelques points du calcul fonctionnel}, Rendiconti del Circolo Matematico di Palermo 22 (1906), 1--74. Metric spaces begin there.
\item Rudin's \emph{Principles of Mathematical Analysis}, 3rd ed., McGraw--Hill, 1976, Chapters~2, 3 and~8, is the model for the metric-space chapters and for the series definition of the exponential. Metric spaces in about forty pages, and the definition of $\exp$, $\cos$, $\sin$ and $\pi$ by power series followed here. The proof of Theorem~\ref{thm:exp}(e) is Rudin's.
\item F.~Q.~Gouv\^ea, \emph{$p$-adic Numbers: An Introduction}, 2nd ed., Springer, 1997. The field $\Q_p$ as a completion, with Hensel's lemma and Ostrowski's theorem.
\item J.~E.~Hutchinson, \emph{Fractals and self-similarity}, Indiana University Mathematics Journal 30 (1981), 713--747. Theorem~\ref{thm:hutchinson}, and a good deal more.
\item M.~F.~Barnsley, \emph{Fractals Everywhere}, Academic Press, 1988. Iterated function systems, the chaos game, and the collage theorem, with pictures.
\item \emph{For the opening stories.} S.~R.~Kaplan, ``The Dottie number'', \emph{Mathematics Magazine} 80 (2007), 73, is where the name and its story appear. The meeting of Steinhaus and Banach in Krak\'ow is told in R.~Ka\l u\.za, \emph{Through a Reporter's Eyes: The Life of Stefan Banach}, Birkh\"auser, 1996. Hausdorff's name for metric spaces is in his \emph{Grundz\"uge der Mengenlehre}, 1914. Hensel introduced the $p$-adic numbers in \emph{\"Uber eine neue Begr\"undung der Theorie der algebraischen Zahlen}, Jahresbericht der Deutschen Mathematiker-Vereinigung 6 (1897), 83--88. Ostrowski's theorem is \emph{\"Uber einige L\"osungen der Funktionalgleichung $\varphi(x)\cdot\varphi(y)=\varphi(xy)$}, Acta Mathematica 41 (1916), 271--284. Gouv\^ea's book proves it.
\end{itemize}
